\documentclass[12pt,a4paper,oneside]{book}

\usepackage[ansinew]{inputenc}
\usepackage[english]{babel}
\usepackage{graphicx}
\usepackage{fourier}
\usepackage{listings}
\usepackage{hyphenat}
\usepackage{xstring}
\usepackage{forloop}
\usepackage{longtable}
\usepackage{enumitem}
\usepackage{array,booktabs}
\usepackage{float}


\lstset{basicstyle=\tt, identifierstyle=\tt, commentstyle=\tt, stringstyle=\tt, keywordstyle=\tt, language=Java, numbers=left, numberstyle=\tiny, numbersep=5pt, tabsize=2, showtabs=false, showspaces=false, showstringspaces=false, prebreak=\mbox{\tiny$\searrow$}, breakindent=0pt, breaklines, breakautoindent=true, captionpos=b}

\oddsidemargin 1,5cm
\textwidth 14,5cm 
 
\begin{document}
\frontmatter

\newsavebox\MyArrowBox%
\sbox\MyArrowBox{$\hookleftarrow$}%
\makeatletter%
\newcommand*{\BreakableChar}{%
  \leavevmode%
  \prw@zbreak%
  %\discretionary{}{}{}%
  %\discretionary{-}{}{}% If want a dash
  \discretionary{\usebox\MyArrowBox}{}{}%
  \prw@zbreak%
}%

\newcounter{index}%
\newcommand{\url}[1]{%
  \texttt{%
  \StrLen{#1 }[\stringLength]%
  \forloop[1]{index}{1}{\value{index}<\stringLength}{%
    \StrChar{#1}{\value{index}}[\currentLetter]%
    {\currentLetter\BreakableChar}%
  }%
  }%
}%

% define here the properties of the title page
\def\title{Erweiterung eines Testdokumentationstools zur besseren/einfacheren Strukturierung und Wartung dokumentierter Testf�lle}
\def\author{Jan Schmid Niederkofler}
\def\advisor{Andrea Janes}
\def\date{Juli, 2017}

\begin{titlepage}
\begin{center}\includegraphics[scale=1.4]{img/logo.pdf}\\\LARGE\vspace{80pt}BACHELOR
THESIS\vspace{50pt}\\{\huge\title}\\\vspace{50pt}\author \end{center}\Large\begin{flushright}\vspace{90pt}Supervisor: \advisor\\\date\end{flushright}
\end{titlepage}

%\chapter*{Abstract}

An abstract should be structured as
follows\footnote{https://www.ece.cmu.edu/~koopman/essays/abstract.html}:

\begin{enumerate}
  \item \textbf{Motivation}: Why do we care about the problem and the results?
  If the problem isn't obviously "interesting" it might be better to put
  motivation first; but if your work is
  incremental progress on a problem that is widely recognized as important, then it is probably better to put the
  problem statement first to indicate which piece of the larger problem you are
  breaking off to work on. This section should include the importance of your
  work, the difficulty of the area, and the impact it might have if successful.

  \item \textbf{Problem statement}:What problem are you trying to solve? What is
  the scope of your work (a generalized approach, or for a specific situation)?
  Be careful not to use too much jargon. In some cases it is appropriate to put
  the problem statement before the motivation, but usually this only works if
  most readers already understand why the problem is important.

  \item \textbf{Approach}: How did you go about solving or making progress on
  the problem? Did you use simulation, analytic models, prototype construction,
  or analysis of field data for an actual product? What was the extent of your
  work (did you look at one application program or a hundred programs in twenty
  different programming languages?) What important variables did you control,
  ignore, or measure?

  \item \textbf{Results}: What's the answer? Specifically, most good computer
  architecture papers conclude that something is so many percent faster,
  cheaper, smaller, or otherwise better than something else. Put the result
  there, in numbers. Avoid vague, hand-waving results such as "very", "small",
  or "significant." If you must be vague, you are only given license to do so
  when you can talk about orders-of-magnitude improvement. There is a tension
  here in that you should not provide numbers that can be easily misinterpreted,
  but on the other hand you don't have room for all the caveats.

  \item \textbf{Conclusions}: What are the implications of your answer? Is it
  going to change the world (unlikely), be a significant "win", be a nice hack,
  or simply serve as a road sign indicating that this path is a waste of time
  (all of the previous results are useful). Are your results general,
  potentially generalizable, or specific to a particular case?
\end{enumerate}

\tableofcontents
\listoftables
\listoffigures
\lstlistoflistings

\mainmatter

\chapter*{Abstract deutsch}

An abstract should be structured as
follows\footnote{https://www.ece.cmu.edu/~koopman/essays/abstract.html}:

\begin{enumerate}
  \item 
 \item \textbf{Problem statement}: Den Verlauf der Testautomatisierung kann man in 2 Phasen unterteilen. Erstere
 die Dokumentation der Testf�lle und letztere das Automatisieren. W�hrend
 die bereits standartisierten Tests aus der Testfactory gewartet werden riskiert
 man durch den st�ndigen Wandel der zu Testenden Programme die Dokumentation zu
 vernachl�ssigen da diese nicht nur Zeitintensiv ist sondern auch nicht so oft
 ben�tigt wird wie ersteres, aber nicht weniger wichtig.
  

  \item \textbf{Motivation}: Die Motivation hinter diesem Projekt ist die Wartbarkeit der Dokumentation zu
  erleichtern bzw. zu erm�glichen. Dessen Funktion vielleicht nicht so viel
  Einfluss auf den Alltag hat, jedoch gro�e Hilfe bieten kann f�r den Fall das
  Fehler entdeckt werden und/oder der zust�ndige f�r den Bereich nicht Anwesend
  ist.
  

  \item \textbf{Approach}:
  Die ersten Schritte bei der Probleml�sung war das Erfassen was bereits
  vorhanden ist womit man Arbeiten kann und was selbst noch zu machen ist.
  Sobald man ein Gef�hl daf�r hat macht man sich mit der Umgebung vertraut, wie
  kann man was Steuern, welche Tools k�nnen helfen, welche Sprache ist hier
  am besten geeignet. Die Wahl fiel schnell auf VBS das es von der Applikation
  am besten unterst�tz wird. Obwohl diese in viele Hinsicht C-Sprachen
  unterlegen w�re.
  Hier ist auch n�tzlich zu Wissen wie weit man das System auslasten kann. Da
  wir hier von tausenden Elementen sprechen m�ssen wir schaugen ab wann wird
  der Ablauf durch die neue Feature langsamer oder es Limitationen vom System aus gibt.
  
  Hier hatten wir noch dazu die M�glichkeit das Tool intern oder extern zu
  erweitern, und entschieden uns auf so viel M�glich auf extern, da sich intern
  bereits �ber die Zeit sich viel Code angesammelt und die externen VBS Scripte
  einfach einzuspielen sind.
  
  Bei der Implementation ist aufzupassen welche Fehler passieren k�nnten und wie
  man mit ihnen Umgehen sollte. Da dieses Projekt gro�es Erweiterungpotential
  hat war es auch wichtig versuchen im Rahmen der Zeit zu bleiben.
  
  

  \item \textbf{Results}:   
  Das Erebnis ist eine leichtere und weniger Fehleranf�llige Eingabe bei der
  Dokumentation. Auch hat man sp�ter eine leichtere Wartung da der Pfad
  standartiesiert eingeben ist und man so Testf�lle ausmachen kann die von der
  �nderung bestimmter Masken betroffen sind. �ltere Systeme m�ssen �berarbeitet
  werden da die Daten in verschiedenen Formaten abgespeichert worden sind.

  \item \textbf{Conclusions}: What are the implications of your answer? Is it
  going to change the world (unlikely), be a significant "win", be a nice hack,
  or simply serve as a road sign indicating that this path is a waste of time
  (all of the previous results are useful). Are your results general,
  potentially generalizable, or specific to a particular case?
  
  Das erreichte Ergebnis bietet die Grundlage einer besserung Wartung. Die
  erleichtert die Eingabe bietet aber sogleich M�glichkeit ausgebaut zu werden
  um mehr Arbeiten weniger Fehleranf�llig sein zu lassen oder gar zu
  automatisieren.
  
\end{enumerate}
\chapter{Einleitung}

\section{Einleitung}
Testf�lle sind die einzelnen Tests die durchgef�hrt werden um sicherzustellen, dass Anwendungen funktionieren wie erwartet.
Die Test�lle werden erstellt von den jeweiligen Verantwortlichen erstellt und
besitzen dann einen fixen Ablauf. Mit HP Application Lifecycle Manager werden diese dokumentiert.
Durch die Erweiterung dieses verbessert sich die Wartbarkeit der dokumentierten Testf�lle. St�ndige Ver�nderungen bei den MAX Anwendungen erschweren diese. Auch verbessert wird die Eingabe der dokumentierten Testf�lle �ber vorgegebenen Dropdown Men�s.

Supply Chain Lufthansa
Lufthansa ist mit fast 126.000 Mitarbeitern eine der gr��ten Fluggesellschaften weltweit. Diese wird aufgeteilt in Passagierbef�rderung, Fracht, Technik, Catering, IT- Dienstleistungen und Service �und Finanzgesellschaft. Dieses Projekt wurde in der Lufthansa Technik realisiert. Genauer im Bereich der Supply Chain. Der Einsatzbereich der Supply Chain ist die Versorgung von Materialen f�r Kunden. Diese werden unterteilt in Material das repariert und wieder im Umlauf gebracht wird, wie ein Flugzeuggetriebe, als auch Material das einfach aufgestockt wird wie Schrauben etc.
Um den Kunden ein Interface bieten zu k�nnen �ber denen sie gew�nschte Bestellungen abgeben k�nnen, bietet Lufthansa verschiedene Software-L�sungen an. Diese unterstehen regelm��igen Updates (bis zu 7 Updates im Jahr)((source)). Um sicherzugehen das nach einem Update alles noch funktioniert wie es sollte wurden bis vor kurzen Integrationstests abgehalten, die daraus bestanden die Applikationen manuell durchzuklicken. Diese Arbeitszeit soll nun mit automatischen Tests reduziert  werden. ((insert source about Lufthansa))
\section{Testdokumentation  / Testautomatisierung}
Die Vorteile der Testautomatisierung sind in erster Linie Zeitersparnis durch die Abnahme der Testf�lle per Computer. Automatisierte Tests laufen, ohne menschliche Interaktion,  in bis zu 20 h durch((insert source, fact check with reality)). Manueller Testf�lle hingegen dauern Tage und bed�rfen Menschen die sie durchklicken. Durch die Reduzierung der menschlichen Interaktion auf ein Minimum entstehen auch weniger Probleme falls ein Experte f�r einen bestimmten Bereich nicht anwesend sein kann, sondern wird nur noch ben�tigt falls die Tests nicht erfolgreich abgeschlossen werden k�nnen. 
Ein Automatisierungszyklus beginnt mit der Auffindung aller Testf�lle. Diese werden dann dokumentiert. Der erste Schritt der Dokumentation beginnt mit einem Tool namens �SAP Workforce Performance Builder�. Dieses ist ein Werkzeug um automatisch bei jeden User Input Screenshots  zu erstellen. So kann ein Testfall in erster Instanz abgebildet sein.
Ein Schritt zur Automatisierung ist die  Dokumentation. Haupts�chlich liegt der Vorteil bei Dokumentation dabei, dass auch wenn ein Experte oder Wissenstr�ger verloren geht, nicht das Wissen selbst verloren ist.
Die Dokumentation beginnt mit �bertragung des Testfalls in das HP Tool �Application Lifecycle Manager� ((insert screenshot)). Hier werden die Schritte einzeln in eine Tabelle eingegeben. Ein einzelner Testschritt besteht aus: wo genau die Aktion stattfindet, was f�r ein Input kommt, welches Objekt an der Maske angesprochen wird. Ein Testschritt kann auch eine SQL abfrage sein oder der 
Start eines Scripts. Diese Dokumentation wird sp�ter verwendet um Testf�lle manuell durchzuspielen. 
Manuelle Tests bedarf es bei Tests die nicht zu Automatisieren sind. Parallel dazu werden manuelle Test verwendet bei fehlgeschlagenen automatisierten Tests den Fehler zu suchen.                                                                                             
 
Ist ein Testfall komplett �bertragen und kontrolliert (4 Augen Prinzip) wird er nach Budapest ((Testfactory/ Testautomatisierung))f reigegeben, wo mit Hilfe von Visual Basic in Script erstellt wird, der den Testfall nachspielt. Dieses Script ist unabh�ngig von der Dokumentation im HPQC, auch wenn dieses darauf basiert. �nderungen an der Maske werden nur in den automatisierten Testf�llen eingespielt. 

Object Action genauer beschreiben + Beschreibungssprache und Formale Sprache (aus dem Wiki) + outsourcing probleme

\section{Ausgangssituation}
Die �bertragung der Testf�lle in HP Quality Center funktioniert, ist durch Textfeld basierten Input aber Fehleranf�llig. Noch dazu ist der Aufwand bei �nderungen zu hoch, so dass diese bei der gro�en Anzahl an Testf�llen die Dokumentation Anpassung  ganz wegbleibt. Die betrifft jetzt nicht automatisierte Tests. Hier das Visual Basic Script angepasst.
((Es funktioniert aber k�nnte besser sein, bei Pflege))
\section{Schwer bis unm�glich manuelle Testf�lle zu warten}
((Anders Beginnen, kein geeigneter erster Unterpunkt, braucht mehr roten Faden))
Ein entscheidender Punkt in der Dokumentation ist die Angabe wo ein bestimmtes Element sich befindet. Dies wird mit Object Path und Object Name angegeben. Bei der �bertragung in das HP tool wird dies manuell eingepflegt. �ndert sich die Maske, durch die regelm��igen Updates, wird diese Angabe falsch. Dies richtig zu Pflegen beinhaltet alle Testf�lle durchzugehen und wo auf die ge�nderte Oberfl�che zugegriffen wird dementsprechend zu �ndern. Folglich sind diese Tests nicht skalierbar und bereits bei wenigen Testf�llen ist mit �nderungen gro�er Aufwand verbunden.

\chapter{Grundlagen und Stand der Technik}
\section{Hp Tool}
Das HP Tool Application Life Cycle Manager, oder kurz HP ALM, ist eine Sammlung
von Software Applikationen die den Lebenszyklus einer Software unterst�tzt.
 In diesen Fall die Erstellung von automatisierten Testf�llen.  
 Spezifisch wird hier HP Quality Center verwendet das eine einzige Plattform f�r
 Test- und Software Zyklus Management bereitstellt. Aufgabe der Software
 Factory ist nicht die genutze Software nach zu bauen sondern eigentst�ndige
 Programme schreiben die auf einer existierend Testumgebung die gew�nschten
 Abl�ufe rekonstruieren. 

\section{Daten}
Eine Sammlung von allen Objekten die bei der Automatisierung angesprochen worden
sind gibt es in Form von XML-Dateien.  Diese liegen extern und nicht im HP-QC
selbst. Da sie von der Testfactory angefertigt worden sind, eignen sie sich
optimal da selbige beim Automatisieren verwendet und somit 
Konflikte mit der Namensgebung der Elemente vermieden werden.

 Da reine XML nicht in dieser Form eingespeichert und verwendet werden k�nnen
 m�ssen diese erst geparst werden. Am besten eigenen sich hier ``Project List''
die von HP bereitgestellt werden. Diese k�nnen auch mit externen Scirpten
erstellt,gel�scht oder bearbeitet werden. Intern sind diese einfach aufrufbar
und somit ideal geeignet um das Eingabefenster je nach Aufruf dynamisch zu
ver�ndern.
\begin{figure}[h]
\includegraphics[width=\textwidth]{\string"img/project lists".png}
\caption{ In diesem Bild sehen wir ein Beispiel von HP eigenen Listen. Diese
k�nnen vom Workflow aus zugegriffen werden und von extern bearbeitet oder
erstellt werden.}
\end{figure}

Der Workflow reguliert den Ablauf der �bertragung der Testf�lle ins HPQC. Mit
diesen kann man Felder an Bedingungen kn�pfen, lehr lassen, ausgrauen,
autovervollst�ndigen etc. 

\section{HP ALM XML}
\begin{lstlisting}
<qtpRep:ObjectRepository xmlns:qtpRep="http://www.mercury.com/qtp/ObjectRepository">
<qtpRep:Objects>
<qtpRep:Object Class="Link" Name="matRequest_info_second_IBO_Materialstamm_Tab">
<qtpRep:ChildObjects>
<qtpRep:Object Class="WinButton" Name="matRequest_info_NotificationBar_FileSave_OpenButton">
\end{lstlisting}

Im oberen Beispiel k�nnen wir die Grundstruktur der XML Dateien erkennen. In der
ersten Zeile wird die Repository URL angegeben. F�r unsere Listen ben�tigen wir
den Unterpunkt ``qtpRep:Objects''. In diesem werden alle Objekte,  von Buttons
�ber Explorer eingeordnet. Die sind unter ``qtpRep:Object'' zu erkennen.

 Die wichtigsten Elemente die wir aus der XML auslesen k�nnen sind die
 Attribute Class und Name. Das Attribut Class gibt mit welche Art von Objekt wir es hier zu
tun haben. Im Code-Beispiel ist es einmal ein Link und darunter ein
Button.
Das n�chste ist f�r uns das wichtigste Element, der Namen mit der ein
Objekte eindeutig gekennzeichnet wird. Am Ende von diesen steht auch eine Art
Klasse in der genau spezifiziert wird was dies darstellt.

Anklickbaren Elemente wie Textboxen etc. werden alle erst untergeordnet in
Explorer, Applikation etc.  unter Childobjects. Ein gr��erer Panel besteht
immer aus mehreren kleinenren Interaktiven Elementen. Diese k�nnen selbst Unterelemente haben.  Mit der
Klasse, Namen und dem Pfad werden sie eindeutig verifiziert und zugeordnet.
werden.

\section{Methodisches Vorgehen}
F�r dieses Projekt benutzen wir die f�r diese Art von Aufgabe untypische
Programmiersprache Visual Basic Script. Aus dem einfachen Grund, dass man mit
dieser einfach auf den Lifecycle Manager zugreifen kann. Auch C(sharp)
bietet diese M�glichkeit ist aber im Funktionsumfang und Ausgereiftheit deutlich
 VBS unterlegen. Auch sind VBS Scripte im HP Alm einfacher einzubinden. 

Noch dazu bietet HP ALM eine einfache Schnittstelle f�r VBS um mit dem Programm
auch von extern kommunizieren zu k�nnen, die Open Test Arcitecture. Der Vorteil
von extern auf die Software zuzugreifen liegt auf dementsprechend externe
Dokumente leichter Zugriff zu haben. In unseren Fall XML-Files. Diese Prozedur
wird unter Implementation genauer erkl�rt.

Ein weitere Vorteil die Scripte nicht direkt in den Hauseigennen ``Worflow''
einzubinden ist das dieser schon mehrere Tausend Zeilen schwer ist. Dies w�rde
die �bersicht des bereits bestehenden Code verringen.

Bevor wir mit der Implementation besch�ftigen, m�ssen wir und vorerst die
Einschr�nkungen unseres Programmes kennen. Unter diesen versteht man
hier die maximale Azahl an Elemente die das Programm erfassen kann, ebenso
die in einer einzigen Iteration uploaden kann oder ob die Ladezeiten dabei
nicht zu gro� werden. Auch d�rfen diese sp�ter nicht die Performance f�r User
beintr�chtigen.

Die Listen die eingetragen werden m�ssen werden auf unter >10.000 gesch�tzt. Die
Tests mit HP ALM zeigen keine negativen Auswirkungen auf HP ALM auch �ber
100.000 Elementen. Auch die Einspielzeit, sofern es in wenigen Iteration
geschieht, bleiben immer unter 15 Minuten. Um bei Listen Elemente hinzuzuf�gen,
wie ein Button etc. muss diese aber bereits im HP QC commited sein, nicht nur
hinzugef�gt. Da aber Listen einzeln commiten mehrere Stunden dauern w�rde muss
man die einzelnen Elemente sp�ter hinzuf�gen.


\chapter{Herausforderungen}

[In der Zusammenfassun so lange nicht aktive Probleme. Dargestellt wie in Zukunft ausbauf�hig und k�nnte und den Umst�nden Probleme machen. Alles komplett abschlie�en, keine Offenen Enden]
\section{Rechte in einen Gro�konzern} 
Bei einen solchen Tool arbeiten zu k�nnen, kann man nicht auf Produktion arbeiten sondern ist es von Vorteil in einer �Sandbox zu arbeiten�. Diese verhindert das beim Arbeiten beim Tool Leute st�rt die gerade mit diesen Arbeiten.
\section{Semaphore / Visual Basic}
Semaphore sind eine Herausforderung da HPQC Zugriff auf die Daten mit semaphoren kann sperren und es an L�sungen gedacht werden muss die immer funktionieren.
\section{XML Parsing}
\section{Zeigen von gro�en Datenmengen auf wenig Feldern}
Die Auswahl von dem gew�nschten Element wird schwierig wenn eine Liste tausende Elemente enth�lt. Eine M�glichkeit dies zu vereinfachen ist die Liste auf mehreren Listen aufzuteielen, die dynamisch generiert werden, abh�ngig von der vorigen Auswahl.
\section{Workflow in HPQC}
Der Workflow in HPQC besteht aus gro�teil unkommentierten 2000 LOC ohne Dokumentation, alles auf nur eine Seite. Dies kann dazu f�hren das leicht die �bersicht verloren geht und Fehler passieren.
\input{chapters/L�sungen}
\chapter*{Problemstelliung}
\chapter{Implementation}
\section{OTA / tdconnect, Visual Basic Script}
Die Open Test Architecture API (OTA) erlaubt es mit tdconnect Zugriff auf HP
Application Lifecycle Manager �ber externe Scripte zu haben. Die Schnittstelle
wird mit Visual Basic Script angesprochen. F�r diesen Zweck wurde Visual Basic
Script C(sharp) bevorzugt da diese Sprache f�r HPQC am geeignetsten ist. VBS ist
gut dokumentiert und verf�gt �ber eine viel gr��ere Bibliothek was die Steuerung des
Tools anbelangt. Der Workflow von HPQC ist auch mit VBS beschrieben, was eine
Andockung mit der selbigen vereinfacht. C(sharp) scheint zu diesen Zeitpunkt
noch nicht f�r die Software ausgearbeitet genug zu sein. 

Creating a connection
\lstset{language=VBScript}

\begin{lstlisting}
' Create a Connection object to connect to Quality Center
   Set tdConnection = CreateObject("TDApiOle80.TDConnection")
'Initialise the Quality center connection
   tdConnection.InitConnectionEx qcURL
'Authenticating with username and password
   tdConnection.Login qcID, qcPWD
'connecting to the domain and project
	tdConnection.Connect qcDomain, qcProject
\end{lstlisting}

�ber dieses tdConnection Objekt k�nnen wir Objekte f�r HP ALM erstellen,
l�schen oder ver�ndern. In unseren Fall sind das selbstgebaute Listen.

\section{ Listen adden + Elemente}
\begin{lstlisting}
' custom is the connection with the elements, oList are the lists access
   Dim custom
    Set custom = tdConnection.Customization
    Dim oLists
    Set oLists = custom.Lists

    'add list for each
    for each item in listToAdd 
                               oLists.RemoveList(lhtListS & item)
                               oLists.AddList(lhtListS & item)
                               'AddList RemoveList; Adding or removing List to the queue on the list object
                next
    custom.Commit
End Sub
\end{lstlisting}


In diesen Beispiel kann man erkennen das die Listen bevor sie geuploadet
werd erst gel�scht werden. Was erst Gegenintuitive erscheinen mag wird so
gehandhabt weil Listen die es bereits gibt einfach �bersprungen werden. Das hat
zur Folge, falls eine Liste nuie Elemente hinzugef�gt bekommt, w�rde diese nicht
aktualisiert werden. Praktisch hier kommt das HP-Alm beim L�schen von Listen den
Befehl einfach �berspringt falls die Liste nicht vorhanden ist. So kommt es zu
keinen Problemen.

\section{Aus XML auslesen}
\begin{lstlisting}
for each child in xmlDoc.SelectNodes("/qtpRep:ObjectRepository/qtpRep:Objects/qtpRep:Object")
  objectLevelName = child.getAttribute("Name")
  objectLevel1.add(filenameXML & "|" & objectLevelName)
  
  'first iteration, getting the second level
  For Each child2 In child.SelectNodes("./qtpRep:ChildObjects/qtpRep:Object")
    objectLevel2Name = child2.getAttribute("Name")
    objectLevel2.add(objectLevelName & "|" & objectLevel2Name )
    
    'third iteration, getting the objects
      For Each child3 In
      child2.SelectNodes("./qtpRep:ChildObjects/qtpRep:Object") objectName =
      child3.getAttribute("Name") objectList.add(objectLevel2Name & "|" & objectName )
      
      'fourth iteration, getting the objects in case objects have objects
        For Each child4 In
          child3.SelectNodes("./qtpRep:ChildObjects/qtpRep:Object") objectName2
         = child4.getAttribute("Name") objectList2.add(objectName & "|" & objectName2 ) 
       next
     next
   next
next
\end{lstlisting}

Das Herzst�ck dieses Projektes, das Durchlaufen der XML. Der erste Schritt ist
das auslesen aller Knotenpunkte mit dem Befehl Select Nodes. Dieser beinhaltet
zuert allte Objekte in einer einzigen Dimension. Wir k�nnen jetzt den ersten
Punkt in die Liste als root Listenname setzen. Dieses Rootelement besteht aus
dem Dateinamen und dem Attributname. Hier gilt zu beachten das maximal 4
Tiefen existieren wobei nur 3 �blich sind und die letze ein Ausnahmefall ist.
Dieser kann bei einer Combobox mit vorgefertigten Elementen vorkommen.

In diesen Fall haben wir eine rekursive linked List erstellt ohne abstrakte
Datentypen, sondern nur �ber den Namen, da VBS diese nicht native unterst�tzt. 
An der linken Seite der Pipe sehen wir wo wir gerade sind, auf der rechten das
gew�nschte Objekt. Der Listenname gibt an an welchen Panel/Browser wir uns
gerade befinden. 
Diese Manipulation ist n�tig da beim einlesen die Level der XML verloren gehen
und alles Eindimensional dargestellt wird und wir nur Knotenpunkt f�r
Knotenpunkt linear verarbeiten k�nnen. 

Die Variante wurde einem rekursiven Aufruf bevorzugt da die Schreibweise
intuitiver erscheint und damit leichter zu warten ist. 

\section{Listnamen Manipulation}
\begin{lstlisting}
for each list in listToAdd
listName = getRightSideOfPipe(list)
for each item in itemToAdd 
  objectList = getLeftSideOfPipe(item)
  if(listName = objectList) Then
    if oLists.IsListExist(lhtListS & list) Then
    Set myList  = oLists.List(CStr(lhtListS & list))'nameList
    Set listRoot = myList.RootNode 
    objectName = getRightSideOfPipe(item)
    On Error Resume Next
      listRoot.AddChild(CStr(objectName))'subname
    On Error GoTo 0
   
    else 
      MsgBox("The List: " + lhtListS & list + " does not exist or is not yet commited") End If
    End If
    next
  next
custom.Commit
\end{lstlisting}

Der Namen ist mit einer Pipe ``|'' getrennt. Auf der rechten Seite zeigt dieser
an welches Objekt gerade im Focus steht. Im Gegensatz auf der linken Seite steht
das Objektlevel, also im welchen Panel, Browser etc das Object platziert ist.
Dieser funktioniert auch als Listenname. Sobald der Listenname nicht mehr mit
dem Objectlevel �bereinstimmt, kreieren wir eine neue Liste.

Der Commit ist der Teil des Programms der am l�ngstem braucht, viele Commits
w�rden viel Zeit beanspruchen, deshalb wird dieser erst ausgef�hrt sobald wir
alle Ebenen und Objekte im Zwischenbuffer oList abgelegt haben.

 *instert graphic with name description

\section{Ergebnis}
\chapter{Future}

\section{Zuk�nftiger Einsatz und Ausbau}
\section{Allen Testf�llen}
Wird bei allen dokumentierten Testf�llen ben�tzt.
[Anstatt dokumentierte Testf�lle -> Dokumentation �Synonyme� -> Standardisieren, einheitliche Beschreibung der Testf�lle, Formalisieren-> sollen alle gleich Aussehen]
\section{ Manuelle Eingabe aus der End-User Maske}
-future
Das Obeject Repository kann nicht nur �ber das Costumization Interface oder Scripten ge�ndert werden sondern auch �ber eine End-User Maske.
\section{Direkt Automatisiert}
-future
Da alle Objekte bereits als Elemente vorhanden sind gibt es die M�glichkeit dies sp�ter zu ben�tzten um zum Teil direkt schon zu automatisieren.
\chapter*{Fazit}


\appendix

\bibliographystyle{plain}                                 
\bibliography{bibliography}

\end{document}